% ============================================================
%  封面与译者说明
% ============================================================
\thispagestyle{plain}

\begin{center}
\vspace*{1.2cm}
{\LARGE\bfseries 曲线的镜像及其 Fukaya 范畴}\\[0.45em]
{\large 三场报告与一份导引}\\[1.0em]
{\large Denis Auroux}\\[1.4em]
{\normalsize
迈阿密同调镜像对称会议（2020 年 1 月 27--29 日）三场报告的手写笔记\\
与 IMSA 成立大会（2019 年 9 月 6 日）报告幻灯片}\\[1.2em]
{\normalsize 中文译注本}

\vspace{0.8em}
\noindent\rule{\textwidth}{0.4pt}
\end{center}

\subsection*{译者说明}
\addcontentsline{toc}{section}{译者说明}

本册合译 Denis Auroux 的两份底本。其一是 2020 年 1 月在迈阿密同调镜像对称会议上所作的三场报告（Talk 1--3）的手写笔记，共 17 页；其二是 2019 年 9 月在 IMSA 成立大会上所作的报告幻灯片，共 14 页，本译作附录 A。两份底本主题连贯：后者是前者的导引，前者是后者在曲线情形下的技术展开。底本两种均无逐字稿，手写笔记只有 LaTeX 化的公式与零散英文，因此译文在保留原有公式编号与记号的前提下，把提纲式的短语补成通顺的句子；凡是补足的内容，均限制在原稿字面允许的范围之内，另在批注中说明。

翻译原则：
\begin{enumerate}
  \item 数学正确性第一。原稿的每一处公式都逐符号核对；凡原稿写法可疑或不完整处，译文照录，另加批注，不改动原稿的数学内容。
  \item 术语先定后用。全书术语见附录 B 术语对照表，该表在开译前冻结，翻译过程中只增不改。
  \item 图全部用 TikZ 重绘。原稿的手绘图按其在页面上的位置与标注重画，图注末尾标出「原稿第 $N$ 页」。原稿中不带标注的速写（例如反复出现的裤子、三价图）只在首次出现处重绘，其后以文字指代。
  \item 不省略证明与计算。原稿中「显然」「同上」之类跳步，凡能以一句话补出者，译文在行文中补出；凡补出之处均以方括号标出。
\end{enumerate}

\textbf{符号约定}：$\K$ 为系数域（底本用 $K$，译文保留 $K$，非阿基米德情形下取 Novikov 域），$T$ 为 Novikov 参数，$\val$ 为赋值，$\Trop$ 与 $\varphi$ 交替表示热带化（原稿两种写法并用，译文不改），$\OO$ 为结构层，$\mathcal L$ 为线丛，$\db$ 为有界导出范畴，$\dsg$ 为奇点范畴，$\Fwr$ 为 wrapped Fukaya 范畴。原稿把 $w_0,w_V,W$ 混用，译文一律照录原稿的写法。

\textbf{批注}：全书共 \textbf{14} 条批注，分「易错」「点拨」「直觉」三类，每条均标明来源，以源文件中的注释形式给出（\texttt{\% [Tutor Note: ...]}），不排入正文，以便与原文对照阅读。

\textbf{底本缺陷}：手写笔记第 9 页之后字号渐大、行距渐疏，第 10、17 页有整页留白；幻灯片第 13 页原有一幅 Jeff Koons 雕塑照片（版权所限，本译本不予复制，正文以文字说明其位置与用意）。第 3 页一处修饰词、第 13 页一处上标记号字迹难辨，均已在疑点清单中标出。

\begin{flushright}
数译 $\cdot$ SY05-28\\
2026 年 11 月
\end{flushright}

\newpage
\tableofcontents
\newpage
